\documentclass[12pt]{article}
\usepackage[utf8]{inputenc}
\usepackage[T1]{fontenc}
\usepackage{hyperref}
\title{Zande Language Resources for K-12 Curriculum Development}
\author{Compiled by Claude Sonnet 5.5 (claude-sonnet-5-5) for Andr\'as Kornai}
\date{}
\begin{document}
\maketitle

\section{Introduction}
Zande (Pazande; ISO 639-3 zne, Glottocode zand1248) is spoken in northeastern DR Congo, the Central African Republic and western South Sudan. It is called Kizande in Lingala. Church sources cited below speak of more than four million Pazande-speaking faithful, a figure that differs from the approximate count in the task description. Zande is among the better-resourced languages in this batch, mainly through religious and literacy materials.

\section{Primary Educational Resources}
\begin{itemize}
\item News report on the Aid to the Church in Need project to teach children to read and pray in their own language in South Sudan, including a Pazande edition of the Child's Bible: \url{https://www.churchinneed.org/south-sudan-teaching-children-to-read-and-pray-in-their-own-language/}
\item Further ACN report on biblical literacy in South Sudan: \url{https://acnuk.org/news/south-sudan-biblical-literacy}
\end{itemize}
No school textbook or primer was found in the search results.

\section{Linguistic and Reference Materials}
\begin{itemize}
\item Wikipedia article on the Zande language: \url{https://en.wikipedia.org/wiki/Zande_language}
\item Wikipedia article on the Zande languages: \url{https://en.wikipedia.org/wiki/Zande_languages}
\end{itemize}
No Zande dictionary or grammar appeared in the search results.

\section{Religious and Community Resources}
\begin{itemize}
\item Vatican News on the first edition of the Pazande Bible in South Sudan (2020): \url{https://www.vaticannews.va/en/africa/news/2020-09/bible-south-sudan-pazande-first-edition.html}
\item Report on translating the Holy Bible into Pazande as a boost to peace: \url{https://encms.abouna.org/node/40190}
\end{itemize}

\section{Digital and Audio Resources}
Global Recordings Network has a Pazande page (Sudan): \url{https://globalrecordings.net/pa/language/79}

\section{Curriculum Development Context}
A complete Pazande Bible and church literacy work provide a standardized written form, which is a useful base for reading materials. Dictionaries, grammars and secular school texts were not found in this search and need further searching. English, Arabic, French and Sango are languages of wider communication across the three countries where Zande is spoken.

\end{document}